\begin{lemma}
Let $Q$ be a completion datum with finite edge and color types, let
$0<\gamma\le\Delta$, $\iota,t>0$, $m=|K|>0$, and let all retained lists
have the same positive cardinality $\ell$. In the product experiment,
for arbitrary real $d,b$,
\[
\mathbb EY_x\le d-t\ \Longrightarrow\
\Pr[Y_x>d]\le e^{-2t^2/m},
\]
\[
\mathbb E|\operatorname{Del}_e|\le(b+\iota/2)\ell
\ \Longrightarrow\
\Pr[|\operatorname{Del}_e|\ge(b+\iota)\ell]
\le \exp\left(-\frac{\iota^2}{8+2\iota}\ell\right).
\]
\end{lemma}
\begin{proof}
By \noderef{NibbleProductSamplingLaw}, the experiment is a probability
measure and every selected edge set is pointwise independent in the line
graph. By \noderef{NibbleSelectedSetProductLaw}, the finite-valued
selected-set maps are measurable and independent, with joint law equal
to the product of their marginal laws. Thus the support predicate for
these values can be taken to be independence in the line graph.

Fix $x$. Take $f(I)=\operatorname{NibbleResidualDegree}(H,I,x)$, as
defined in \noderef{NibbleResidualDegree}. For any two families of
independent edge sets differing only at one color,
\noderef{NibbleResidualDegreeBoundedDifference} gives
$|f(I)-f(J)|\le1$. Apply \noderef{FiniteIndependentBoundedDifferences}
to this function, the selected-set maps, and the support predicate just
specified. It yields $\Pr[Y_x\ge\mathbb EY_x+t]\le e^{-2t^2/m}$.
The mean hypothesis implies $\mathbb EY_x+t\le d$, so
$\{Y_x>d\}\subseteq\{Y_x\ge\mathbb EY_x+t\}$, giving the first bound.

Fix $e$. For each retained color $a$, put
$R_a(S)\iff\exists f\in S$ adjacent to $e$ in the line graph.
By \noderef{NibbleDeletedColors}, the number of successes of these
predicates among $a\in M(e)$ is exactly $Z=|\operatorname{Del}_e|$.
Apply \noderef{FiniteIndependentBernoulliCountTail} with this retained
color set of size $\ell$ and increment $s=\iota\ell/2>0$. This gives
\[
\Pr[Z\ge\mathbb EZ+s]\le\exp(-s^2/(2\mathbb EZ+s))
\le\exp(-s^2/(2\ell+s)).
\]
The helper includes the zero-mean case and supplies $\mathbb EZ\le\ell$;
both denominators are positive, which justifies the second inequality.
Our mean hypothesis gives
$\mathbb EZ+s\le(b+\iota)\ell$, so the event in the statement is
contained in this non-strict tail event. Finally, since $\ell>0$,
\[
\frac{(\iota\ell/2)^2}{2\ell+\iota\ell/2}
=\frac{\iota^2\ell}{8+2\iota}.
\]
Both helper applications are to the original product measure: their
finite push-forward proofs transport the integrals and event measures,
so no change of experiment or unstated identification of means is needed.
\end{proof}
